\documentclass[10pt,a4paper]{amsart}
\usepackage[margin=19mm]{geometry}
\usepackage{amsmath,amssymb,mathtools}
\usepackage[scr=rsfs,cal=euler]{mathalfa}
\usepackage{xfrac}
\usepackage[unicode,hidelinks,bookmarks=false]{hyperref}
\newcommand{\ie}{i.e.\ }
\newcommand{\cf}{cf.\ }
\newcommand{\resp}{respectively\ }
\newcommand{\Subsection}{\subsection}
\providecommand{\nobreakdash}{\nobreak\hbox{-}}
\setlength{\emergencystretch}{2em}
\multlinegap0pt
\usepackage{mathtools}
\newtheorem{lemma}{Lemma}
\newcommand{\R}{\mathbb R}
\newcommand{\T}{\mathbb T}
\title{Dimension-free estimates for full discrete maximal functions associated with Euclidean balls and spheres}
\author{Mahdi Hormozi, Jakub Niksi\'nski,, B{\l}a\.zej Wr\'obel}
\date{Source: arXiv:2609.10763v1 (CC BY 4.0)}
\begin{document}
\maketitle

\noindent\emph{Source and licence.} Dimension-free estimates for full discrete maximal functions associated with Euclidean balls and spheres. Version arXiv:2609.10763v1 (CC BY 4.0), distributed by arXiv under the Creative Commons Attribution 4.0 International licence, http://creativecommons.org/licenses/by/4.0/, as recorded in the arXiv licence metadata for this exact version and shown on the abstract page https://arxiv.org/abs/2609.10763v1. Full licence notice: \texttt{licenses/arxiv-2609.10763-CC-BY-4.0.txt}.\par\smallskip

\noindent\textit{Editorial scope.} Counted unit: the article's lemma lem:euclidean-quadratic-phase (source0.tex lines 1018--1030). The article's complete original proof of it is delivered with it (source0.tex lines 1032--1168, one \texttt{proof} environment of 137 lines). No mathematics is written, repaired or re-derived: the statement and the proof are the article's own lines, in the article's own notation, and no retained line is substituted or re-flowed.  No further source-local context is needed: every cross-reference of every retained line resolves inside the two delivered spans.  Nothing was omitted from any retained span, so the package carries no omission record.  No retained line contains a citation, so the wrapper carries no bibliography and no bibliography entry.  No retained line contains a figure environment, an image-inclusion command or drawing code, so no image is delivered and the package declares no asset dependency.  Editorial formatting only, in addition: an AMS \texttt{amsart} preamble that restates the article's own declarations and macros used by the retained lines, namely the environments \texttt{lemma} on separate counters, together with the standard packages those lines need.  The retained environments print in this document as Lemma 1 in order of appearance, whereas the article prints the same items with the numbering of its own full document.  The complete native source of the article as distributed in the arXiv e-print for this exact version is bundled unchanged as \texttt{sources/209947/source0.tex}.
\begin{lemma}\label{lem:euclidean-quadratic-phase}
There is an absolute constant $c>0$ such that, for every integer $L\geq2$ and every
$t,\theta\in\R$ we have
\begin{equation}\label{eq:euclidean-finite-rigidity}
 1-|Q_L(t,\theta)|^2
 \geq c\min\left\{1,
 \min_{\varepsilon\in\{0,1\}}
 \left(
 L^4\|t-\varepsilon/2\|_{\T}^2
 +L^2\|\theta-\varepsilon/2\|_{\T}^2
 \right)\right\}.
\end{equation}
\end{lemma}

\begin{proof}
First we will prove that for every integer $M\geq2$ and every $a\in\mathbb \R$, one has
\begin{equation}\label{eq:euclidean-linear-phase}
 1-\left|\frac1M\sum_{j=0}^{M-1}e(ja)\right|^2
 \geq c\min\{1,M^2\|a\|_{\T}^2\}.
\end{equation}
We may assume that $a \in (-1/2, 1/2]$. Let $\varepsilon_0>0$ be a fixed small constant, to be determined later and let 
\begin{equation*}
    D_M(a)= \frac{1}{M} \sum_{j=0}^{M-1} e(ja)
\end{equation*}
denote the normalized Dirichlet kernel. 

In the case $M|a| \le \varepsilon_0$ we use the formula for Dirichlet's kernel and Taylor's theorem to get 
\begin{align*}
1-\left|D_M(a)\right|^2
&=1-\frac{\sin^2(\pi Ma)}{M^2\sin^2(\pi a)}=1-
\frac{\Big(\pi Ma-\frac{\pi^3M^3a^3}{6}
+O(M^5|a|^5)\Big)^2}
{M^2\Big(\pi a-\frac{\pi^3a^3}{6}
+O(|a|^5)\Big)^2}\\
&=1-
\frac{\Big(1-\frac{\pi^2M^2a^2}{6}
+O(M^4|a|^4)\Big)^2}
{\Big(1-\frac{\pi^2a^2}{6}
+O(|a|^4)\Big)^2}=1-
\frac{1-\frac{\pi^2M^2a^2}{3}
+O(M^4|a|^4)}
{1-\frac{\pi^2a^2}{3}
+O(|a|^4)}\\
&=\frac{\pi^2(M^2-1)a^2}{3}
+O(M^4a^4)=\frac{\pi^2M^2|a|^2}{3}
\left(1-\frac{1}{M^2}+O(\varepsilon_0^2)\right).
\end{align*}
Now, as long as $\varepsilon_0$ is small enough, the last expression is bounded from below by $cM^2|a|^2$ for some absolute constant $c>0$.

If $M|a| \in (\varepsilon_0, 1)$, then we use the fact that the function 
\[
 [0,1] \ni x \mapsto \frac{\sin(\pi x)}{x}
\]
is decreasing. This way, since $M \ge 2$, one obtains
\begin{align*}
    \left|D_M(a)\right|&= \left| \frac{\sin(\pi Ma)}{M \sin(\pi a)}  \right|= \frac{\sin(\pi M|a|)/(M|a|)}{ \sin(\pi|a|)/(|a|)} 
    \\
    &\le \frac{\sin(\pi M|a|)/(M|a|)}{ \sin(\pi M |a|/2)/(M|a|/2)}  = \cos(\pi M|a|/2) \le \cos(\pi\varepsilon_0/2) \le 1-c_0
\end{align*}
for some constant $c_0>0$. Lastly, if $M|a| \ge 1$, then using the inequality $\sin (\pi x) \ge 2x$, which holds for $x \in [0, 1/2]$, we get
\begin{align*}
    \left|D_M(a)\right|&= \left| \frac{\sin(\pi Ma)}{M \sin(\pi a)}  \right|= \frac{|\sin(\pi Ma)|}{ M\sin(\pi|a|)} \le  \frac{1}{2M |a|} \le \frac{1}{2}.
\end{align*}
In conclusion, after combining all the cases \eqref{eq:euclidean-linear-phase} is proved.
\par
Now we turn to the proof of \eqref{eq:euclidean-finite-rigidity}.
Let $\delta_0>0$ be a small constant to be determined later and let
\begin{equation*}
    \delta= 1- \left|Q_{L}(t, \theta) \right|^2, \qquad N=2L+1.
\end{equation*}
If $\delta>\delta_0$, then \eqref{eq:euclidean-finite-rigidity} holds trivially, hence we consider only the case $\delta \le \delta_0$.
First we will show that if $\delta_0$ is small enough, then 
\begin{equation} \label{eq: bound on t}
    \min_{\varepsilon \in \{0,1\}}\|t-\varepsilon/2 \|_{\T} \le C \frac{\sqrt{\delta}}{L^2}
\end{equation}
for some absolute constant $C>0$.

Expanding the square and changing variable $h=k-l$ in the spirit of Weyl we get 
\begin{equation} 
\label{eq: Weyl diff}
\begin{split}
    |Q_{L}(t, \theta)|^2&= \frac{1}{(2L+1)^2} \sum_{k,l=-L}^L e\big(t(k^2-l^2)+\theta(k-l) \big) \\
    & =\frac{1}{(2L+1)^2} \sum_{h=-2L}^{2L} e\big( th^2+\theta h\big) \sum_{l=\max(-L,-L-h)}^{\min(L,L-h)}e(2t hl).
\end{split}
\end{equation}
Using \eqref{eq: Weyl diff} and \eqref{eq:euclidean-linear-phase} one obtains
\begin{align*}
    \delta&=1-|Q_{L,\theta}(t)|^2 \ge  \frac{1}{N^2} \sum_{h=-2L}^{2L} \Big(N-|h| \Big) \Big(1- \left|D_{N-|h|}(2 ht) \right|\Big) \\
    &\ge \frac{1}{2N^2} \sum_{h=-2L}^{2L} \Big(N-|h| \Big) \Big(1- \left|D_{N-|h|}(2 ht) \right|^2\Big) \\
    &\ge  \frac{c}{L} \sum_{h=\lceil L/2 \rceil}^L  \min\{1, L^2 \|2t h \|_{\T}^2 \}.
\end{align*}

From the above we conclude that for sufficiently large constant $C>0$ the proportion of $h \in \{ \lceil L/2 \rceil, \ldots, L \}$ satisfying
\begin{equation*} 
    \| 2th \|_{\T} \le \frac{C\sqrt{\delta}}{L}
\end{equation*}
is greater than $3/4$, in particular there are two consecutive elements $h_0,h_0+1$ of that interval satisfying the inequality above. Then by triangle inequality we get 
\begin{equation*}
    \|2 t\|_{\T} \le  \|2 t(h_0+1)\|_{\T} +  \|2 th_0\|_{\T} \le \frac{2C\sqrt{\delta}}{L}.
\end{equation*}
Thus for some $\varepsilon \in \{0,1\}, m \in \mathbb Z, \eta \in \mathbb R$ we have 
\begin{equation*}
    t=  m+\varepsilon/2+ \eta, \quad \left |\eta \right| \le  \frac{C\sqrt{\delta}}{L}.
\end{equation*}
For $h_0$ as before we have 
\begin{equation*}
    \|2 \eta h_0 \|_{\T}= \|2th_0 \|_{\T} \le \frac{C\sqrt{\delta}}{L},
\end{equation*}
moreover $2 |\eta|h_0 \le 2C\sqrt{\delta_0}$, so if $\delta_0$ is small enough we get $|2 \eta h_0|<1/2$ and hence from previous inequality we see that
\begin{equation*}
    2|\eta|h_0 \le \frac{C\sqrt{\delta}}{L}.
\end{equation*}
Since $h_0 \ge L/2$, the inequality above implies
\begin{equation} \label{eq: t-eps bound}
    \|t-\varepsilon/2 \|_{\T}=|\eta| \le \frac{C'\sqrt{\delta}}{L^2},
\end{equation}
which establishes \eqref{eq: bound on t}. 

Now we will show that for the same choice of $\varepsilon$ as for $t$, we have 
\begin{equation*} 
%\label{theta bound}
    \| \theta -\varepsilon/2 \|_{\T} \le \frac{C'' \sqrt{\delta}}{L}.
\end{equation*}
Using the fact that $|Q_{L}(t,\theta)|=|Q_{L}(t+ 1/2,\theta+1/2)|$ one can see that 
$|Q_L(t, \theta)|=|Q_L(\eta, \beta)|$ for some $\beta \in [-1/2, 1/2]$ such that 
\begin{equation*}
    |\beta|= \|\theta -\varepsilon/2 \|_{\T}.
\end{equation*}
Note that we have 
\begin{equation}
\label{eq: eta, beta}
\begin{split}
    \delta=1-|Q_L(\eta,\beta)|^2&=\frac{1}{N^2} \sum_{k,l=-L}^{L}\Big(1-\cos\big(2\pi\eta(k^2-l^2)+2\pi\beta(k-l)\big) \Big) \\
    &\gtrsim \frac{1}{N^2} \sum_{k,l=-L}^{L} \left\| \eta(k^2-l^2)+\beta(k-l)\right \|_{\T}^2, 
\end{split}
\end{equation}
where at the end we have used the inequality $1- \cos(2\pi x) \gtrsim  \|x \|_{\T}^2$. Applying \eqref{eq: t-eps bound} and \eqref{eq: eta, beta} we obtain 
\begin{align*}
    \frac{1}{N^2} \sum_{k,l=-L}^{L} \left\| \beta(k-l)\right \|_{\T}^2 &\lesssim \frac{1}{N^2} \sum_{k,l=-L}^{L} \left\| \eta(k^2-l^2)+\beta(k-l)\right \|_{\T}^2 \\
    &+  \frac{1}{N^2} \sum_{k,l=-L}^{L}\left\| \eta(k^2-l^2) \right \|_{\T}^2 \lesssim \delta + L^4 |\eta|^2 \lesssim \delta.
\end{align*}
However, note that due to \eqref{eq:euclidean-linear-phase} and the comparison $1- \cos(2 \pi x) \lesssim \|x \|_{\T}^2$, which holds for all $x \in \mathbb R$, we get
\begin{align*}
    c\min\{1, L^2\|\beta \|_{\T}^2 \} \le1-|D_N(\beta)|^2&=  \frac{1}{N^2} \sum_{k,l=-L}^{L} \big(1-\cos\left(2\pi\beta(k-l) \right) \big)   \\ &\lesssim\frac{1}{N^2} \sum_{k,l=-L}^{L} \left\| \beta(k-l)\right \|_{\T}^2 \lesssim \delta,
\end{align*}
the above combined with definition of $\beta$ and \eqref{eq: t-eps bound} gives 
\begin{equation*}
    L^4\|t - \varepsilon/2 \|_{\T}^2+ L^2 \| \theta - \varepsilon/2 \|_{\T}^2 \lesssim \delta
\end{equation*}
for some $\varepsilon \in \{0,1\}$, which concludes the proof.
\end{proof}

\end{document}
